\documentclass{handout}
\usepackage{handoutSetup}\usepackage{handoutShortcuts}  

%\renewcommand{\cite}{\citeyearpar}
%\newboolean{AnswersTF}
\setboolean{AnswersTF}{true}


\begin{document}
\handoutHeader

\begin{verbatimwrite}{\jobname.title}
The Blanchard (1985) Model of Perpetual Youth
\end{verbatimwrite}

\centerline{\LARGE \input \jobname.title}
\medskip 
%\centerline{\url{http://www.econ2.jhu.edu/people/ccarroll/public/lecturenotes}}
\medskip\medskip

This handout analyzes a way to relax the standard assumption of infinite
lifetimes in the Ramsey/Cass-Koopmans
growth model.  The trick, introduced by
Blanchard~\citeyearpar{blanchardFinite}, is to assume that the economy
is populated by agents who face a constant probability of death.
Thus, an agent who has lived a thousand years is no more likely to die
in the next year than an agent who was born yesterday.

Time is measured continuously.  If there is an instantaneous
probability $\pDies$ of dying, the probability (as viewed by a person
alive in period $t$) of still being alive (not dead) in period
$\tThen$ is
\begin{equation}\begin{gathered}\begin{aligned}
        \Alive_{t}^{\tThen} & =  e^{-\pDies(\tThen-t)} . \label{eq:living}
\end{aligned}\end{gathered}\end{equation}

\begin{enumerate}
\item Assuming the instantaneous utility function is $\uFunc(c) = 
\log c$, and that the pure rate of time preference is $\timeRate$, 
explain why the objective function at time $t$ of an individual 
household in this model will be to maximize
\begin{displaymath}
        \int_{t}^{\infty} (\log c_{\tThen}) e^{-(\timeRate+\pDies)(\tThen-t)} d\tThen. \label{eq:objfunc}
\end{displaymath}
where $c_{\tThen}$ is the consumer's consumption at time $\tThen$.  For convenience,
you may wish to define
\begin{equation}\begin{gathered}\begin{aligned}
        \hat{\timeRate} & =  \timeRate+\pDies. \label{eq:thetahat}
\end{aligned}\end{gathered}\end{equation}

\ifthenelse{\boolean{AnswersTF}}{

\begin{quote}
{\it Answer:}

A consumer with a pure time preference rate of zero will downweight 
the utility he receives conditional on being alive by the probability 
that he is still alive, yielding a discounted utility of
\begin{equation}\begin{gathered}\begin{aligned}
        \int_{t}^{\infty} (\log c_{\tThen}) \Alive_{t}^{\tThen} e^{-\timeRate (\tThen-t)}d\tThen & =  \int_{t}^{\infty} (\log c_{\tThen}) e^{-\timeRate (\tThen-t)} e^{-\pDies (\tThen-t)} \nonumber d\tThen \\
         & =    \int_{t}^{\infty} (\log c_{\tThen}) e^{-(\timeRate+\pDies) (\tThen-t)} d\tThen
\\       & =    \int_{t}^{\infty} (\log c_{\tThen}) e^{-\hat{\timeRate} (\tThen-t)} d\tThen
\end{aligned}\end{gathered}\end{equation}
where $\hat{\timeRate} = \timeRate+\pDies$.

A similar point holds in the discrete time model.  A sensible thing to
assume is that if you have died before $t+1$, you get zero utility in
$t+1$ and after.  Thus, in the two-period context, if the probability
of death between $T-1$ and $T$ was zero we would have
\begin{displaymath}
   V_{T-1} = \max \uFunc(C_{T-1}) + \left(\frac{1}{1+\timeRate}\right) \uFunc(C_{T})
\end{displaymath}
while if there is a probability $\pDies$ of dying between $T$ and $T+1$
value would be:
\begin{equation}\begin{gathered}\begin{aligned}
   V_{T-1} & =  \max \uFunc(C_{T-1}) + (1-\pDies)\left(\frac{1}{1+\timeRate}\right) \uFunc(C_{T}) + \pDies*\beta*0 \nonumber 
\\   & =  \max \uFunc(C_{T-1}) + \left(\frac{1-\pDies}{1+\timeRate}\right) \uFunc(C_{T}) \nonumber
\\   & \approx  \max \uFunc(C_{T-1}) + \left(\frac{1}{1+\timeRate+\pDies}\right) \uFunc(C_{T}) \label{eq:approx} 
\end{aligned}\end{gathered}\end{equation}

But behavior in this case is virtually indistinguishable from the 
behavior that would be induced if the consumer had a time preference 
rate of $\timeRate+\pDies$.  In continuous time, the approximation in 
\eqref{eq:approx} becomes exact.


\end{quote}
}{\pagebreak

}

If the probability of death is constant, the expected remaining life 
for an agent of any age is given by
\begin{displaymath}
        \int_{0}^{\infty} \pDies \tau e^{-\pDies \tau}d\tau = \pDies^{-1}.
\end{displaymath}
which we will call the agent's `horizon.'  For example, if the chances 
of dying per year are $\pDies=1/50$, then the agent's horizon is 50 years.

We will assume that at every instant of time, a large cohort, whose size
is normalized to be $\pDies$, is born.  

\item For an economy that has existed forever, explain why the 
formula for the aggregate population at time $t$, $P_{t}$, will be
\begin{equation}\begin{gathered}\begin{aligned}
        P_{t} & =  \pDies \int_{-\infty}^{t} \Alive_{s}^{\tNow} ds \label{eq:totpop}
\end{aligned}\end{gathered}\end{equation}
and show that this formula implies that the population is $P_{t}=1$.  

\ifthenelse{\boolean{AnswersTF}}{

\begin{quote}
{\it Answer:}

The aggregate population will be the sum of the still-alive persons 
from all past generations.

The proportion of a population born at time $s$ that is still living 
at time $\tThen$ is $\Alive_{s}^{\tThen}$ by equation \eqref{eq:living}.  
Thus, the absolute size at time $\tThen$ of a cohort of size $\pDies$ born at 
time $s$ is $\pDies \Alive_{s}^{\tThen} = \pDies e^{-\pDies (\tThen-s)}$.  The economy's 
total population will be the sum of the populations of all the cohorts 
that are currently living.  Since the economy has existed forever, 
there will be remaining members of every cohort back to $s=-\infty$.  
Thus, indexing each cohort by a time index $s$, \eqref{eq:totpop} is 
simply the sum of the populations of all currently living members of 
every generation.

Substituting the formula for $\Alive$, the integral becomes
\begin{equation}\begin{gathered}\begin{aligned}
        P_{t} & =  \pDies \int_{-\infty}^{t} e^{-\pDies(t-s)} ds 
\\        & =  \pDies \int_{t}^{\infty} e^{-\pDies(s-t)} ds 
\\        & =  \pDies \int_{0}^{\infty} e^{-\pDies \tThen}d\tThen
\\        & =  \pDies \pDies^{-1}  \label{eq:fmhint}
\\        & =  1 
\end{aligned}\end{gathered}\end{equation}
where \eqref{eq:changevars} comes from a change of variables $\tThen = s-t$ 
and \eqref{eq:fmhint} follows from the hint.

\end{quote}
}{\pagebreak

}

Now some notation.  We will define variable $x(s,t)$ as the value at 
date $t$ of the variable $x$ for a consumer who was born at date $s$.  
Thus, $c(s,t)$ is consumption at $t$ of a consumer born at $s$.

Suppose that the consumers in this economy do not have a bequest 
motive.  If they have positive assets at the instant when they die, 
they are no happier than if they had zero assets.  This means that if 
someone were willing to pay them something while they are still alive 
for the right to inherit their assets whenever they die, these consumers would 
happily take that deal.  

For a consumer with wealth $w(s,t)$ who has probability of dying $\pDies$, 
the flow value of the right to inherit that wealth is $\pDies w(s,t)$.  We 
will therefore assume that insurance companies exist that pay a 
consumer with wealth $w(s,t)$ an amount $\pDies w(s,t)$ in exchange for the 
right to receive that consumer's wealth when he dies.  (The insurance 
company will make zero profits).

But notice that from the standpoint of the consumer, this is 
equivalent to saying that the interest rate received on wealth is 
higher by amount $\pDies$.

Now suppose the marginal product of capital in this 
perfectly-competitive economy is constant at $\rfree$ and suppose an agent 
born in $s$ receives exogenous labor income in period $t$ of $y(s,t)$.  
This plus the insurance scheme implies that the agent's dynamic budget 
constraint is given by
\begin{equation}\begin{gathered}\begin{aligned}
        \dot{w}(s,t) & =  (\rfree+\pDies) w(s,t) + y(s,t)-c(s,t) \label{eq:budconstr}.
\end{aligned}\end{gathered}\end{equation}

Define the `effective' interest rate as viewed by a consumer as $\hat{\rfree}  = \rfree+\pDies$.

\item Write the current-value Hamiltonian for the 
individual's maximization problem and use it to show that the growth 
rate of consumption in period $t$ for a consumer born at time $s$ is 
given by
\begin{equation}\begin{gathered}\begin{aligned}
        \left(\frac{\dot{c}(s,t)}{c(s,t)}\right) & =  \hat{\rfree} - \hat{\timeRate}
\\  & =  \rfree-\timeRate.     
\end{aligned}\end{gathered}\end{equation}

\ifthenelse{\boolean{AnswersTF}}{

\begin{quote}
{\it Answer:}

The current-value Hamiltonian is written
\begin{displaymath}
        \Ham(c,w,\lambda) = \log c(s,t) + \lambda(\hat{\rfree} w(s,t) + y(s,t) - c(s,t))
\end{displaymath}
so the first optimality condition $\partial \Ham/\partial c(s,t) = 0$ implies
\begin{equation}\begin{gathered}\begin{aligned}
        1/c(s,t) & =  \lambda
\\  - \dot{c}(s,t)/c(s,t)^{2} & =  \dot{\lambda}       
\end{aligned}\end{gathered}\end{equation}
and the second optimization condtion implies
\begin{equation}\begin{gathered}\begin{aligned}
        \dot{\lambda} & =  \hat{\timeRate} \lambda - \partial \Ham/\partial w(s,t)  
\\       & =  \hat{\timeRate} \lambda - \hat{\rfree} \lambda  
\\      \left(\frac{\dot{\lambda}}{\lambda}\right) & =  \hat{\timeRate}-\hat{\rfree}
\\  \left(\frac{-\dot{c}(s,t) c(s,t)}{c(s,t)^{2}}\right) & =  \hat{\timeRate}-\hat{\rfree}
\\  \left(\frac{\dot{c}(s,t)}{c(s,t)}\right) & =  \hat{\rfree}-\hat{\timeRate} 
\\   & =  \rfree-\timeRate     .
\end{aligned}\end{gathered}\end{equation}

\end{quote}
}{\pagebreak

.
\pagebreak
}


\item Use the first order condition for consumption and the 
intertemporal budget constraint implied by \eqref{eq:budconstr} to 
show that the level of consumption in time $t$ for an individual born 
at time $s$ is
\begin{displaymath}
        c(s,t) = \hat{\timeRate}(w(s,t)+h(s,t))
\end{displaymath}
where recall that $\hat{\timeRate}=\timeRate+\pDies$ and human wealth is
\begin{displaymath}
        h(s,t) = \int_{t}^{\infty} \left(y(s,\tThen)/\hat{\mathcal{R}}_{t}^{\tThen}\right) d\tThen
\end{displaymath}
where
\begin{equation}\begin{gathered}\begin{aligned}
  \hat{\mathcal{R}}_{t}^{\tThen} & =  e^{\int_{t}^{\tThen} (\rfree_{\mu}+\pDies) d\mu}
\end{aligned}\end{gathered}\end{equation}
(this is simply the compound discount factor necessary to take
account of time-varying interest rates; if interest rates are
constant at $\rfree$ it reduces to the usual term $e^{-\rfree(\tThen-t)}$).

\ifthenelse{\boolean{AnswersTF}}{

\begin{quote}
{\it Answer:}

The IBC says that the PDV of consumption must equal wealth plus the PDV
of future labor income:
\begin{equation}\begin{gathered}\begin{aligned}
        \int_{t}^{\infty} (c(s,\tThen)/\hat{\mathcal{R}}_{t}^{\tThen})  d\tThen & =  w(s,t) + \int_{t}^{\infty} y(\tThen,t)/\hat{\mathcal{R}}_{t}^{\tThen} d\tThen
\\      \int_{t}^{\infty} (c(s,\tThen)/\hat{\mathcal{R}}_{t}^{\tThen}) d\tThen & =  w(s,t) + h(s,t).
\end{aligned}\end{gathered}\end{equation}

But if $\dot{c}(s,t)/c(s,t) = \hat{\rfree}_{t}-\hat{\timeRate}$ then
\begin{equation}\begin{gathered}\begin{aligned}
        c(s,\tThen) & =  c(s,t)e^{\int_{s}^{\tThen}\hat{\rfree}_\mu d\mu}e^{-\hat{\timeRate}(\tThen-t)}
\end{aligned}\end{gathered}\end{equation}
implying
\begin{equation}\begin{gathered}\begin{aligned}
        \int_{t}^{\infty} (c(s,\tThen)/\hat{\mathcal{R}}_{t}^{\tThen}) d\tThen & =  
\int_{t}^{\infty} c(s,t) e^{\int_{s}^{\tThen}\hat{\rfree}_\mu  d\mu}e^{-\hat{\timeRate}(\tThen-t)}e^{-\int_{s}^{\tThen}\hat{\rfree}_\mu d\mu} d\tThen \nonumber
\\         & =  c(s,t) \int_{t}^{\infty} e^{-\hat{\timeRate}(\tThen-t)} d\tThen
\\   & =  c(s,t)/\hat{\timeRate}
\end{aligned}\end{gathered}\end{equation}
so that the IBC becomes:
\begin{equation}\begin{gathered}\begin{aligned}
        c(s,t)/\hat{\timeRate} & =  w(s,t)+h(s,t)  \\
        c(s,t) & =  \hat{\timeRate} (w(s,t)+h(s,t)) .
\end{aligned}\end{gathered}\end{equation}

\end{quote}
}{\pagebreak

}

Suppose we define upper-case variables as the aggregate value across all
generations currently living of the corresponding lower-case value, e.g.
aggregate consumption is
\begin{equation}\begin{gathered}\begin{aligned}
        C(t) & =  \int_{-\infty}^{t} \pDies c(\tThen,t) \Alive_{\tThen}^{t} d\tThen.
\end{aligned}\end{gathered}\end{equation}

\item (Easy) Show that the aggregate level of consumption in 
this economy is
\begin{equation}\begin{gathered}\begin{aligned}
        C(t) & =  \hat{\timeRate} (W(t) + H(t)) . \label{eq:cagg}
\end{aligned}\end{gathered}\end{equation}

\ifthenelse{\boolean{AnswersTF}}{

\begin{quote}
{\it Answer:}
\begin{equation}\begin{gathered}\begin{aligned}
        C(t) & =   \int_{-\infty}^{t} \pDies c(\tThen,t) \Alive_{\tThen}^{t} d\tThen \\
         & =  \int_{-\infty}^{t} \pDies \hat{\timeRate}(w(\tThen,t)+h(\tThen,t)) \Alive_{\tThen}^{t} d\tThen \\
         & =  \hat{\timeRate} (W(t)+H(t))
\end{aligned}\end{gathered}\end{equation}
from the definition of $W(t)$ and $H(t)$.

\end{quote}
}{\pagebreak

}

Suppose all living agents in this economy receive the same noncapital 
income, $y(s,t) = Y(t)$.  Since every member of the population has the 
same income, and the size of the population is one, aggregate income 
will also be $Y(t)$.  Aggregate human wealth at $t$ is therefore
\begin{equation}\begin{gathered}\begin{aligned}
        H(t) & =  \int_{t}^{\infty} (Y(\tThen)/\hat{\mathcal{R}}_{t}^{\tThen}) d\tThen
\end{aligned}\end{gathered}\end{equation}
where 
\begin{equation}\begin{gathered}\begin{aligned}
        \dot{H}(t) & =  \hat{\rfree}_{t} H(t) - Y(t) .
\end{aligned}\end{gathered}\end{equation}

The differential equation for aggregate wealth can be 
shown to be
\begin{equation}\begin{gathered}\begin{aligned}
        \dot{W}(t) & =  w(t,t) - \pDies W(t) + \int_{-\infty}^{t} \dot{w}(\tThen,t) \pDies e^{-\pDies(t-\tThen)} d\tThen, \label{eq:Wdot}
\end{aligned}\end{gathered}\end{equation}
where $w(t,t)=0$ is the wealth of newly born generations, $\pDies W(t)$ is 
the wealth of those who are dying at the moment, and the last term is 
the change in wealth for those who neither die nor are born in this 
period.  But \eqref{eq:budconstr} implies that
\begin{equation}\begin{gathered}\begin{aligned}
  \int_{-\infty}^{t} \dot{w}(\tThen,t) \pDies e^{-\pDies(t-\tThen)} d\tThen & =  (\rfree+\pDies) W(t) + Y(t)-C(t)
\end{aligned}\end{gathered}\end{equation}
so \eqref{eq:Wdot} becomes
\begin{equation}\begin{gathered}\begin{aligned}
        \dot{W}(t) & =  \rfree W(t) + Y(t)-C(t).
\end{aligned}\end{gathered}\end{equation}

Collecting, writing out $\hat{\rfree}=\rfree+\pDies$ and $\hat{\timeRate}=\timeRate+\pDies$, and 
dropping the $(t)$ arguments gives us the following equations for 
aggregate variables:
\begin{equation}\begin{gathered}\begin{aligned}
        C & =  (\pDies+\timeRate) (H+W) \label{eq:clevel} \\
        \dot{H} & =  (\rfree+\pDies) H - Y  \\ 
        \dot{W} & =  \rfree W + Y - C %\label{eq:wdot}
\end{aligned}\end{gathered}\end{equation}

\input BlanchardFiniteQ.tex


\end{enumerate}


\ifthenelse{\boolean{AnswersTF}}
{

\pagebreak

 
\input handoutBibMake
}{}

\end{document}
